\section{第一条线：从Hall快重联到“什么才是必要条件”}\label{story:fast}
这一线索先问一个有限的问题：已经启动的薄层能否在弱耗散条件下快速处理磁通？Biskamp的双流体计算和GEM比较给出肯定的模型证据；Rogers进一步把快率联系到色散波和开放排气。后来的研究把这个解释变成可以检验的必要性命题，进而发现没有快速色散波的快支、以及电子各向异性改变几何的另一条路径。因而这些论文的关系是“建立基准—提出机制—寻找反例—收紧结论”，不是后一篇简单宣布前一篇失效。

阅读时必须保留两层机制：Hall或压力响应怎样组织排气，电子冻结又由哪一个非理想项打破。第二层不因宏观速率对某耗散系数较弱敏感就可以省略。以下把模型开关、扫描路径和实际破冻结项并排检查，再在下一条线追问上游场和层长由谁决定。

\subsection{Hall解耦与快率：GEM解决了什么，没解决什么}
\citet{Biskamp1997}在准中性、均匀密度、标量压力的双流体模型中发现，电子层随电子尺度变化，而某些宏观通量率对电子质量和电子黏性弱敏感。计算包含正则化耗散，实际伪谱计算还采用六阶扩散。电子Kelvin--Helmholtz湍流主要关联有限耗散率，并没有显著增加原本已快的重联率。

该文表II按Alfvén时间归一化后，$d_i$ 从0.4降到0.05，峰值率由约0.68降到0.34；作者明确没有确定 $d_i\to0$ 极限。固定密度也移除了某些压力动力学。因此不能把它改写为“任意系统尺寸、无耗散且普适0.1”。

\citet{Birn2001}的GEM比较使用有限磁岛种子，刻意跳过线性触发阶段。Hall、混合与PIC模型的非线性通量曲线较接近，峰值约0.24；然而粒子模型中的电子电流层和离子喷流宽度可比对应流体模型大约两倍。文章还指出过量耗散可降低重联率，并把结论限定在二维。

因此GEM提供的是\textbf{同一已启动算例的部分宏观量可跨模型复现}的证据。它并未证明自然触发、压力张量、电子层厚度、能量分配和三维不稳定性均可由一个Hall项替代。

\subsection{色散波与开放排气：有力的模型内关联，有限的普适性}
\citet{Rogers2001}把whistler和kinetic Alfvén波的二次色散，与电子流出随层厚变化及排气开口联系起来。其等温两流体扫描中，两类分支都允许时重联率约0.12；仅一个分支时约0.06；强导引场、冷等离子体且两类二次色散分支均被抑制时，电流层较长，速率显著下降。

控制参数包括投影Alfvén速度、声速和质量比，而不是单独的 $B_g/B_{\rm rec}$。将线性分类用于非线性重联时，相关重联磁场分量自身随时间和位置变化；作者用上游反向场作方便的参数替代，并承认这一点。这不是适用于任意电子压力闭合的必要条件定理。下文电子CGL机制说明，同样的开放排气还应检查压力项而非仅凭波模名称解释。

2015年的两篇强导引场研究进一步给出\emph{同类两流体框架内}的反例：在$d_e\gtrsim\rho_s$、没有快速色散波的参数区，仍可得到对耗散弱敏感的快重联，并以PIC作对照。它们还重新检查2001年的扫描路径，发现启动期上游场耗尽与初始片宽也参与决定测得速率；详见本章第\ref{sec:withoutfdw}节。因而本节的早期结果应读作特定扫描中的关联，不能继续写成快重联的必要条件。\citep{L216,L217}

\subsection{没有快速色散波的快重联：反例具体推翻了什么}\label{sec:withoutfdw}
\citet{L216}在强导引场两场模型中分别研究$d_e>\rho_s$和$\rho_s>d_e$。前者没有通常用来解释快重联的快速色散分支。超电阻从$10^{-11}$增到$10^{-8}$时，两流体算例的归一化率仍在约0.025--0.029；对应单流体率则近于$\eta_H^{0.252}$。这是同一类耗散扫描中的实测差别，不是从X形图像直接判定快慢。

其关键新发现是\textbf{微观耗散层的上游场和长度会自调节}。当耗散层比$h=\max(d_e,\rho_s)$薄，上游$B_\delta$随厚度一起减小；某些区间内层长度也随厚度缩短。因此，对固定$B_\delta$和固定内层长度推导出的耗散依赖，不一定是实际算例采取的路径。第一章式\eqref{eq:two_layer_rate}把这一观察重建成两层匹配。它解释给定$B_h,w_v$时的弱耗散依赖，却没有自洽选出这两个外部量。

长文\citet{L217}进一步比较流体与PIC。在有快速色散波的一组参数中，两模型峰值率相差不到约10\%，但启动和饱和时刻不同；PIC早期出现的次级岛、喷流弯曲也没有在流体中逐一重现。在无快速色散波的\emph{小盒}算例中，没有次级岛仍能快速重联。因此，该组证据同时排除了“快速色散波必须存在”和“次级岛必须出现”这两种过强的必要性表述；并未排除它们在其他参数区影响速率。

\subsubsection{重做早期扫描：区分改变声尺度和改变电子尺度}
长文第VII节重做类似Rogers等的扫描：固定$d_e$而减小$\rho_s$时，速率下降可以出现；但固定$\rho_s$、在仍无快速色散波的区间增大$d_e$，速率也会上升。两组都改变了$h$相对于初始电流片宽的比例，因而改变启动等待时间及启动时还能供应多强的上游场。按当时测得的$B_h^2$归一化后，差别明显缩小。这里的结论是早期扫描不能单独隔离色散波的因果作用，并非“线性波对重联毫无作用”。\citep{Rogers2001,L217}

\subsubsection{尺寸扫描没有同时放大所有方向}
Harris系列把流出方向长度从基准的0.5倍变到8倍，基准$L_{x0}=100h$；入流方向盒宽与初始片厚保持不变。局地上游场取离中面约$4h$的位置，按入流穿越时间平均，所得流体与PIC率约0.09--0.17。长文还排除最后约四分之一时窗，以减少可用磁通耗尽的影响。这个结果支持所测试几何中的弱尺寸依赖，但不能替代全几何自相似放大，也不能证明任何边界下都为0.1。

学习时可先对照短文图1、3，再读长文图14--19：前者检查耗散层如何变化，后者检查尺寸、归一化和扫描路径。两篇包含相关算例，不能作为两次完全独立的统计验证。\citep{L216,L217}

\subsection{电子CGL的新机制：快率与四极场极性反转}
\citet{Cassak2015}逐项关闭Hall、电子惯性，或改用离子CGL、仅把电子各向异性放入动量方程，构成机制对照。特定电子CGL进入欧姆定律时，仍可得到约0.06--0.1的快率和开放排气；仅离子CGL或仅动量各向异性则较慢。面外四极扰动的极性与通常Hall四极场相反，是比“存在四极图样”更具体的判据。

不过模拟始终保留电阻0.005和四阶耗散，压力方程还省略电磁功加热；因此快几何机制不等于CGL压力在X点独自打破冻结。原文提到另一电子状态方程的慢结果只来自私人通信，PIC相符也未给出对照图。这些内容应作为进一步检验方向，不能与已展示的消融证据同权。

\subsection{早期闭合与移除Hall对照：不能只比较开关}
\citet{L029}的Hall/混合比较采用特定经验电子压力张量闭合，且其方程省略磁旋转项；这与今天完整十矩不是同一模型。
它提示离子非对角压力输运能改变层结构，不能从全局磁通相近推断离子应力无关。
\citet{L154}对无Hall快速重联的检查更具体：粗格Sweet--Parker层约2.5个网格时出现岛，约5格时可恢复较稳定层；去掉Hall后的率仍约0.015，含Hall约0.03，但仍有超电阻等耗散。
因此该对照揭示分辨率与耗散的作用，不能单独证明所有系统中Hall既必要又充分。

\citet{L035}的边界压缩算例中，混合模型因离子回旋效应出现更宽的响应，进入快速阶段的时间与Hall模型不同；局部嵌入Hall区可复现其特定总通量演化。
这解决的是\emph{该驱动条件下局地两流体区域是否足以改变整体响应}，未解决任意装置边界下的模型耦合。

\paragraph{这一线索走到了哪里。}这条线已建立的可靠结论是：多个不同微观实现能在特定已启动算例中产生相近的局地快率；快速色散波、四极图样、开放X形中的任一个，都不是单独判定快重联的充分证据。下一步的判别应保持外部供磁、初始片宽和电子非理想机制可比，分别扫描微观尺度，测量瞬时上游场后再归一。尚未被选出的长度和供磁条件，正是下一节的起点。
